% ============================================================================
% 第五章：总结与展望
% 优化说明：精简结构，内容深入，避免复杂嵌套
% ============================================================================

\section{总结与展望}

{\footnotesize

% ============================================================================
% 研究总结
% ============================================================================
\subsection{研究总结}

\begin{frame}{研究背景与核心问题}
    \begin{block}{MI-BCI 应用前景}
        \begin{itemize}
            \item \textbf{康复医疗}：脑卒中、脊髓损伤患者重建运动功能
            \item \textbf{人机交互}：为严重运动障碍患者提供新型控制方式
            \item \textbf{神经工程}：深化对大脑运动皮层 ERD/S 现象的理解
        \end{itemize}
    \end{block}

    \begin{block}{关键挑战：时间窗选择}
        \begin{itemize}
            \item \textbf{ERD 动态特性}：事件相关去同步化存在显著个体间差异（潜伏期 0.5--2s）
            \item \textbf{EEG 信号特性}：非平稳、低信噪比（SNR 约 $-10$ 至 $-20$ dB）
            \item \textbf{实时性要求}：在线 BCI 系统需毫秒级响应（延迟 $< 100$ms）
            \item \textbf{资源约束}：嵌入式设备内存通常 $< 1$MB，功耗 $< 100$mW
        \end{itemize}
    \end{block}

    \begin{block}{核心研究问题}
        如何在保持分类性能的前提下，实现轻量化、可解释、自适应的时间窗优化？
    \end{block}
\end{frame}

\begin{frame}{现有方法局限性分析}
    \begin{table}[h]
        \centering
        \renewcommand{\arraystretch}{1.3}
        \begin{tabular}{l|l|l}
            \hline
            \textbf{方法} & \textbf{局限性} & \textbf{复杂度} \\
            \hline
            固定时间窗 & 无法适应个体间差异 & $O(1)$ \\
            滑动窗 & 计算复杂度线性增长 & $O(n)$ \\
            贝叶斯优化 & 高维参数空间效率低 & $O(n^3)$ \\
            深度强化学习 & 百万参数、训练不稳定、黑盒决策 & $O(|\theta|)$ \\
            \hline
        \end{tabular}
    \end{table}

    \vspace{0.3cm}
    \begin{block}{核心矛盾}
        \begin{itemize}
            \item \textbf{DQN 的需求}：百万级参数、大量训练数据（$>10^5$ 样本）、GPU 加速
            \item \textbf{BCI 的现实}：数据稀缺（典型实验 $<1000$ 试次）、设备资源受限
            \item \textbf{根本问题}：模型复杂度与任务难度不匹配
        \end{itemize}
    \end{block}
\end{frame}

\begin{frame}{研究方法：表格型 Q-Learning}
    \begin{block}{马尔可夫决策过程形式化}
        \begin{itemize}
            \item \textbf{状态空间} $\mathcal{S}$：136 个离散时间窗状态（起始点 0--0.8s，终点 0.8--3.0s，步长 0.2s）
            \item \textbf{动作空间} $\mathcal{A}$：4 个边界调整动作 $\{t_s^-, t_s^+, t_e^-, t_e^+\}$
            \item \textbf{奖励函数} $R(s,a)$：CSP-SVM 5 折交叉验证准确率
            \item \textbf{状态转移} $P(s'|s,a)$：确定性过程
            \item \textbf{折扣因子} $\gamma = 0.95$：平衡即时与长期奖励
        \end{itemize}
    \end{block}

    \begin{block}{算法设计}
        \begin{itemize}
            \item \textbf{表格型表示}：维护 $Q(s,a)$ 表，参数数量 $=|\mathcal{S}|\times|\mathcal{A}|=544$
            \item \textbf{更新规则}：$Q(s,a) \leftarrow Q(s,a) + \alpha[r + \gamma\max_{a'}Q(s',a') - Q(s,a)]$
        \end{itemize}
    \end{block}

    \begin{block}{核心优势}
        轻量化（2KB 内存）$\cdot$ 收敛性保证（Watkins, 1989）$\cdot$ 可解释性强
    \end{block}
\end{frame}

\begin{frame}{实验结果与性能对比}
    \begin{center}
    \renewcommand{\arraystretch}{1.4}
    \begin{tabular}{lcc}
        \hline
        \textbf{指标} & \textbf{Table Q} & \textbf{DQN} \\
        \hline
        平均准确率 & 62.68\% & 62.70\% \\
        标准差 & $\pm 8.12\%$ & $\pm 8.05\%$ \\
        统计检验 & \multicolumn{2}{c}{$t(8) = 0.024, p = 0.981$} \\
        效应量 (Cohen's $d$) & \multicolumn{2}{c}{$0.003$（可忽略）} \\
        \hline

    \end{tabular}
    \end{center}

    \begin{block}{核心结论}
        \begin{itemize}
            \item 表格型 Q-Learning 在保持与 DQN 相当性能的同时，显著降低计算资源需求
            \item 性能差异无统计学意义 ($p > 0.05$)，效应量接近零
        \end{itemize}
    \end{block}
\end{frame}

\begin{frame}{计算效率详细对比}
    \begin{table}[h]
        \centering
        \renewcommand{\arraystretch}{1.3}
        \begin{tabular}{lccc}
            \hline
            \textbf{指标} & \textbf{Table Q} & \textbf{DQN} & \textbf{提升倍数} \\
            \hline
            参数量 & 544 & 1,211,776 & 2227$\times$ \\
            内存占用 & 2 KB & 250 KB & 125$\times$ \\
            推理延迟 & 0.001 ms & 1.2 ms & 1200$\times$ \\
            训练时间/被试 & 2.5 min & 7.8 min & 3.1$\times$ \\
            FLOPs/推理 & $\sim 10^2$ & $\sim 10^6$ & 10000$\times$ \\
            \hline
        \end{tabular}
    \end{table}
    
    \vspace{0.3cm}
    \begin{block}{嵌入式部署意义}
        \begin{itemize}
            \item Table Q 可在 ARM Cortex-M4（主频 80MHz, SRAM 320KB）上实时运行
            \item DQN 需高性能处理器（如 Cortex-A7）或专用 AI 加速器
            \item 功耗对比：Table Q $< 10$mW vs DQN $> 500$mW
        \end{itemize}
    \end{block}
\end{frame}

\begin{frame}{主要理论贡献}
    \begin{block}{轻量化 BCI 设计新范式}
        \begin{itemize}
            \item \textbf{核心思想}：将深度强化学习简化为表格型强化学习
            \item \textbf{设计原则}：模型复杂度与任务难度匹配，避免过度参数化
            \item \textbf{实现效果}：轻量化与高性能的统一，参数量减少 2227 倍
            \item \textbf{推广价值}：适用于其他状态空间有限的 BCI 优化任务
        \end{itemize}
    \end{block}

    \vspace{0.3cm}
    \begin{block}{算法选择准则}
        \begin{itemize}
            \item \textbf{适用边界}：状态空间 $|\mathcal{S}| < 10^4$ 且状态 - 动作对可充分访问
            \item \textbf{表格型优势}：收敛性保证、样本效率高、可解释性强、无超参数调优
            \item \textbf{DQN 适用场景}：$|\mathcal{S}| > 10^6$ 或连续状态空间
        \end{itemize}
    \end{block}
\end{frame}

\begin{frame}{主要实践贡献}
    \begin{block}{可解释性 BCI 系统实践}
        \begin{itemize}
            \item \textbf{Q 值表可视化}：可直接展示每个状态的最优动作和预期回报
            \item \textbf{学习轨迹追踪}：记录 Q 值收敛过程，分析智能体学习动态
            \item \textbf{决策透明度}：医生/研究者可以理解系统为何选择特定时间窗
            \item \textbf{临床转化价值}：可解释性是医疗 AI 系统审批的关键前提
        \end{itemize}
    \end{block}

    \vspace{0.3cm}
    \begin{block}{实验验证体系}
        \begin{itemize}
            \item \textbf{对比方法}：10 种系统性评估（固定窗、DQN、4 种 Q-Learning 变体等）
            \item \textbf{统计检验}：配对样本 $t$ 检验 + 效应量分析 (Cohen's $d$)
            \item \textbf{被试分析}：9 名被试个体级别性能详细对比
            \item \textbf{消融实验}：验证各算法组件的贡献
        \end{itemize}
    \end{block}
\end{frame}

% ============================================================================
% 研究局限性
% ============================================================================
\subsection{研究局限性}

\begin{frame}{研究局限性 1/3}
    \begin{block}{验证范围的局限}
        \begin{itemize}
            \item \textbf{数据集单一}：仅 BCI Competition IV 2a，缺乏跨数据集验证
            \item \textbf{被试群体局限}：9 名健康被试，未包含脑卒中患者
            \begin{itemize}
                \item 患者脑电信号非平稳性更强、信噪比更低
                \item 可能存在病灶侧运动皮层功能重组
            \end{itemize}
            \item \textbf{验证场景}：离线分析，未在线闭环验证
            \item \textbf{统计效力}：样本量 $N=9$，统计检验效力约 0.6
        \end{itemize}
    \end{block}

    \vspace{0.3cm}
    \begin{block}{方法设计的局限}
        \begin{itemize}
            \item \textbf{离散化步长}：0.2 秒，可能无法定位全局最优的连续时间点
            \item \textbf{间接奖励}：通过 CSP-SVM 准确率间接构建，非直接衡量时间窗质量
            \item \textbf{干扰因素}：分类准确率受 CSP 特征质量、SVM 参数等影响
            \item \textbf{潜在改进}：探索 ERD 强度、信噪比、互信息等直接奖励信号
        \end{itemize}
    \end{block}
\end{frame}

\begin{frame}{研究局限性 2/3}
    \begin{block}{Q-Learning 变体研究不足}
        \begin{itemize}
            \item \textbf{已对比变体}：Standard、Double、Dueling、Dueling Double Q-Learning
            \item \textbf{现象观察}：Double Q-Learning 性能略低于 Standard Q-Learning
            \item \textbf{机制分析不足}：
            \begin{itemize}
                \item 双 Q 网络的过估计修正在本任务中可能过于保守
                \item 状态空间有限、奖励确定性高时，Standard Q-Learning 过估计问题不显著
            \end{itemize}
            \item \textbf{未探索变体}：Prioritized Experience Replay、Noisy Nets、Distributional RL
            \item \textbf{Q 值表演化分析缺失}：未分析收敛速度、Q 值分布时空模式
        \end{itemize}
    \end{block}
\end{frame}

\begin{frame}{研究局限性 3/3}
    \begin{block}{通用性的局限}
        \begin{itemize}
            \item \textbf{分类器依赖}：仅 CSP-SVM，未验证 EEGNet、Shallow FBCSP 等深度学习模型
            \item \textbf{BCI 范式单一}：仅运动想象，未验证 SSVEP、P300 等范式
            \item \textbf{嵌入式测试缺乏}：未在 ARM Cortex-M 系列单片机上部署验证
        \end{itemize}
    \end{block}

    \vspace{0.3cm}
    \begin{block}{分析维度的局限}
        \begin{itemize}
            \item \textbf{理论分析为主}：参数计数、FLOPs、复杂度分析
            \item \textbf{缺乏实证基准}：未测量实际运行时间、功耗、温度影响
        \end{itemize}
    \end{block}
\end{frame}

% ============================================================================
% 未来工作展望
% ============================================================================
\subsection{未来工作展望}

\begin{frame}{未来工作 1/3：在线和临床验证}
    \begin{block}{在线实时验证}
        \begin{itemize}
            \item \textbf{系统移植}：将 Q 表移植到 ARM Cortex-M4/M7 嵌入式系统
            \item \textbf{评估指标}：
                \item 端到端延迟：从 EEG 采集到反馈呈现（目标 $< 100$ms）
                \item 在线准确率：实时反馈场景下的分类性能
                \item 信息传输率 (ITR)：bits/min，综合评估系统性能
                \item 用户体验：NASA-TLX 认知负荷量表、系统可用性量表 (SUS)
        \end{itemize}
    \end{block}

    \begin{block}{临床验证}
        \begin{itemize}
            \item \textbf{目标群体}：脑卒中后运动功能障碍患者、脊髓损伤患者
            \item \textbf{试验设计}：样本量 $N \geq 30$，4--8 周干预，每周 3--5 次训练
            \item \textbf{评估工具}：Fugl-Meyer 运动功能评分、Barthel 指数
            \item \textbf{预期价值}：为实际应用价值提供循证医学证据
        \end{itemize}
    \end{block}
\end{frame}

\begin{frame}{未来工作 2/3：算法效率优化}
    \begin{block}{快速奖励估计方法}
        \begin{itemize}
            \item \textbf{基于元学习}：学习奖励函数的代理模型，新被试只需少量试次
            \item \textbf{缓存机制}：复用历史被试交叉验证结果，基于相似度检索
            \item \textbf{目标}：离线训练时间从分钟级缩短到秒级，实现"即插即用"
        \end{itemize}
    \end{block}

    \begin{block}{时间窗精度提升}
        \begin{itemize}
            \item \textbf{更精细离散化}：探索 0.1 秒步长(状态数增至约 500，仍远小于 DQN)
            \item \textbf{混合优化策略}：第一阶段：表格型 Q-Learning 粗搜索；第二阶段：贝叶斯优化/梯度方法精搜索
            \item \textbf{连续时间建模}：探索连续时间强化学习方法
        \end{itemize}
    \end{block}
\end{frame}

\begin{frame}{未来工作 3/3：泛化与迁移}
    \begin{block}{算法泛化能力}
        \begin{itemize}
            \item \textbf{深度学习结合}：与 EEGNet、Shallow FBCSP 等模型结合
            \item \textbf{空间层面优化}：
            \begin{itemize}
                \item 通道选择：基于强化学习选择最优 EEG 通道子集
                \item CSP 空间滤波器优化：联合优化分量数和参数
                \item 时空联合优化：统一建模为 MDP，实现时空特征协同优化
            \end{itemize}
        \end{itemize}
    \end{block}

    \vspace{0.3cm}
    \begin{block}{迁移学习策略}
        \begin{itemize}
            \item \textbf{核心思路}：从不同被试的最优 Q 表中提取通用先验知识
            \item \textbf{技术路线}：聚类分析 $\rightarrow$ 原型提取 $\rightarrow$ 快速适配
            \item \textbf{预期收益}：校准时间从 30 分钟缩短至 5 分钟，降低使用门槛
        \end{itemize}
    \end{block}
\end{frame}

% ============================================================================
% 本章小结
% ============================================================================
\subsection{本章小结}

\begin{frame}{本章小结}
    \begin{block}{研究总结}
        \begin{itemize}
            \item \textbf{核心方法}：提出表格型 Q-Learning 轻量化时间窗优化方法
            \item \textbf{性能表现}：与 DQN 无显著差异 ($t(8)=0.024, p=0.981, d=0.003$)
            \item \textbf{四大贡献}：
            \begin{itemize}
                \item 轻量化 BCI 设计新范式
                \item 算法选择准则（$|\mathcal{S}| < 10^4$ 时表格型优于 DQN）
                \item 可解释性 BCI 系统实践
                \item 完善的实验验证体系
            \end{itemize}
        \end{itemize}
    \end{block}

    \begin{block}{研究局限性}
        验证范围~$\cdot$ 方法设计~$\cdot$ 变体研究~$\cdot$ 通用性~$\cdot$ 分析维度
    \end{block}

    \begin{block}{未来方向}
        在线临床验证~$\cdot$ 算法效率优化~$\cdot$ 泛化能力探索~$\cdot$ 迁移学习
    \end{block}
\end{frame}

\begin{frame}{研究的理论意义与实践价值}
    \begin{block}{理论意义}
        \begin{itemize}
            \item \textbf{强化学习在 BCI 中的应用边界}：明确表格型与深度强化学习的适用场景
            \item \textbf{轻量化 AI 设计原则}：为资源受限场景下的模型选择提供理论指导
            \item \textbf{可解释性 BCI 方法论}：提供可解释性系统设计的实践参考
        \end{itemize}
    \end{block}

    \begin{block}{实践价值}
        \begin{itemize}
            \item \textbf{便携式 BCI 设备}：使高性能时间窗优化可在嵌入式设备部署
            \item \textbf{康复医疗应用}：为脑卒中患者提供可负担的 BCI 康复方案
            \item \textbf{技术转化潜力}：低功耗、低成本、可解释性支撑临床审批
        \end{itemize}
    \end{block}

    \begin{block}{结语}
        \centering
        \large
        轻量化设计~$\cdot$ 性能保证~$\cdot$ 可解释性~$\cdot$ 临床转化
    \end{block}
\end{frame}

} % end of \footnotesize
